% --- Translation notes ---
% -------------------------
%
% The introduction promises to name the build-or-reveal move *every time* one is
% made. That makes it a recurring structural element of the book, not incidental
% prose, so it gets a first-class environment: retrofitting one across 121
% sections later would be miserable, and prose-only notes drift in wording.
%
% Built from rules and spacing only. Upstream loads none of framed, mdframed,
% tcolorbox or etoolbox, and adding a package to the pinned submodule's
% dependency set for a box border is a poor trade.

\newlength{\translationNoteSkip}
\setlength{\translationNoteSkip}{0.6\baselineskip}

% \begin{translationNote}{Build from scarcity}
%   Python's dict is this table, finished. We build it anyway, because ...
% \end{translationNote}
\newenvironment{translationNote}[1]{%
  \par%
  \needspace{4\baselineskip}%
  \vspace{\translationNoteSkip}%
  \noindent{\color{MediumGray}\rule{\linewidth}{0.4pt}}\par%
  \vspace{-0.5\baselineskip}%
  \noindent{\footnotesize\sffamily\bfseries\color{DeepGray}#1}\par%
  \nopagebreak%
  \vspace{0.2\baselineskip}%
  \begingroup%
  \small%
  \setlength{\parindent}{0pt}%
  \ignorespaces%
}{%
  \par%
  \endgroup%
  \vspace{-0.3\baselineskip}%
  \noindent{\color{MediumGray}\rule{\linewidth}{0.4pt}}\par%
  \vspace{\translationNoteSkip}%
}

% The two moves the introduction commits to. Wrappers rather than free text so
% the vocabulary stays fixed across the whole book and the notes stay greppable.
\newenvironment{buildFromScarcity}{%
  \begin{translationNote}{Build from scarcity}%
}{%
  \end{translationNote}%
}
\newenvironment{revealTheSurface}{%
  \begin{translationNote}{Reveal the surface}%
}{%
  \end{translationNote}%
}

% Points where the two languages genuinely diverge, as opposed to points where
% we have made a choice between them. These recur throughout -- they are the
% seam the introduction calls the most interesting boundary in the book -- so
% they get their own environment rather than being folded into prose.
\newenvironment{departure}{%
  \begin{translationNote}{Where Python departs from Lisp}%
}{%
  \end{translationNote}%
}

% For the third category the introduction insists on keeping distinct from the
% other two: a claim that is a structural fact about the language, not a matter
% of taste.
\newenvironment{notTaste}{%
  \begin{translationNote}{Fact, not taste}%
}{%
  \end{translationNote}%
}
